\section{电压与线路电流的安全裕度计算}\label{sec:method}

本节给出安全接入容量算法的实际计算过程。对每一组逆变器下垂增益，首先在选定负荷水平下求解非线性 DistFlow，得到各节点电压、线路有功及无功潮流。然后考虑新增负荷和光伏功率在指定范围内的变化，计算电压和电流可能达到的边界值，最后与设备运行限值比较。对于整个负荷区间均满足约束的配置，算法继续增加接入容量。

\subsection{节点电压响应与分块 Hessian}\label{sec:hess}

取新增负荷区间 $I=[\lambda_0-h,\lambda_0+h]$，以 $\lambda_0$ 为参考负荷，$s=\lambda-\lambda_0$ 为负荷变化量，$w$ 为光伏有功和无功扰动。对节点 $i$，平方电压随负荷和光伏变化的局部表达式为
\begin{equation}
\nu_i(\lambda_0+s,w)
=\nu_i^0+\frac{\partial\nu_i}{\partial\lambda}s+
(\nabla_w\nu_i)^\top w+\varepsilon_i(s,w),
\label{eq:voltage_local}
\end{equation}
其中 $\nu_i^0$ 是参考工况的潮流结果，两个一阶导数分别描述\emph{增加负荷}和\emph{光伏出力变化}对节点电压的影响，$\varepsilon_i$ 表示非线性余项。

计算上述系数时，将非线性潮流方程写为 $F(z,\lambda,\bm g,w)=0$，$z=(\nu,\ell)$。数值潮流求解器首先获得参考状态及雅可比矩阵 $J=F_z$；由同一雅可比矩阵求解线性方程组得到一阶灵敏度，再求二阶灵敏度：
\begin{align}
z_t&=-J^{-1}F_t,\label{eq:zt}\\
z_{t_at_b}&=-J^{-1}D^2F[(z_{t_a},e_a),(z_{t_b},e_b)] .
\label{eq:ztt}
\end{align}
其中 $t=(\lambda,w)$。实际实现采用线性方程组求解，不需要显式计算 $J^{-1}$。一阶灵敏度用于预测电压变化趋势，二阶灵敏度用于估计负荷区间和光伏扰动内的预测误差。

为直接判断电压是否满足 $V_{\min}\le V_i\le V_{\max}$，分别计算下限和上限对应的裕度违反量
$\phi_i^-=V_{\min}^2-\nu_i$ 和 $\phi_i^+=\nu_i-V_{\max}^2$。
对任一 $\phi$，将二阶误差分为三项，其区间最大值分别记为
\begin{equation}
|\phi_{\lambda\lambda}|\le B_{\lambda\lambda},\qquad
\norm{\nabla_w\phi_\lambda}_2\le B_{\lambda w},\qquad
\norm{\nabla_w^2\phi}_2\le B_{ww}.
\label{eq:blocks}
\end{equation}
这三个系数对应三个直接的工程影响：$B_{\lambda\lambda}$ 描述新增负荷引起的电压响应非线性，$B_{\lambda w}$ 描述负荷与光伏同时变化的耦合，$B_{ww}$ 描述光伏扰动自身引起的非线性。各系数从相应节点的二阶灵敏度中提取，用于计算当前负荷单元的电压变化范围。

\subsection{光伏扰动下的节点电压范围}

本算例中，第 $j$ 台光伏的扰动满足式~\eqref{eq:U} 的半圆约束。对于电压一阶变化中的系数对 $(a,b)$，对应半圆扰动能够产生的最大贡献可直接计算为
\begin{equation}
\supp_{\U_j}(a,b)=
C_j\begin{cases}
\sqrt{a^2+b^2},&a\ge0,\\
|b|,&a<0.
\end{cases}
\label{eq:support}
\end{equation}
因此，在计算每个节点的最不利电压时，可以直接使用其有功和无功灵敏度，而无需逐个枚举光伏注入组合。

设 $R_w=(\sum_jC_j^2)^{1/2}$。对节点电压违反量 $\phi$，根据一阶灵敏度和上一小节的三个二阶系数，计算整个光伏扰动集合上的上界
\begin{align}
\Phi_\phi(s)={}&
\phi_0+\phi_\lambda s+
\sum_{j=1}^{m}
\supp_{\U_j}(\phi_{w_{p,j}},\phi_{w_{q,j}})
\nonumber\\
&+\tfrac12B_{\lambda\lambda}s^2+
B_{\lambda w}R_w|s|+
\tfrac12B_{ww}R_w^2.
\label{eq:Vbound}
\end{align}
式中前两项由参考潮流和负荷灵敏度计算，第三项为光伏扰动的最大一阶影响，最后三项是非线性修正。将负荷单元的两个端点 $s=\pm h$ 分别代入上式，即可得到电压约束的单元检查量：
\begin{equation}
\max\{\Phi_{\phi_i^-}(-h),\Phi_{\phi_i^-}(h),
\Phi_{\phi_i^+}(-h),\Phi_{\phi_i^+}(h)\}\le0.
\label{eq:voltage_accept}
\end{equation}
对每个节点重复这一步，便得到该负荷单元的电压安全判定。

\subsection{线路有功、无功潮流与电流计算}

线路是否过载取决于线路电流。在非线性 DistFlow 中，支路 $l$ 的接收端有功和无功潮流分别记为 $P_l=a_l$、$Q_l=b_l$，接收端电压幅值为 $V_l=\sqrt{\nu_l}$，于是
\begin{equation}
I_l=\frac{\sqrt{P_l^2+Q_l^2}}{V_l}.
\label{eq:current_physical}
\end{equation}
工程计算中，我们先利用潮流求解器得到参考工况的 $P_l,Q_l,V_l$，再用灵敏度计算新增负荷和光伏扰动对 $P_l,Q_l$ 的影响，最终计算线路可能达到的最大电流。

具体地，将某条线路的有功、无功潮流组合成向量 $y_l=(P_l,Q_l)^\top$，其在当前负荷单元内的计算表达式为
\begin{equation}
y_l(s,w)=y_l^0+t_ls+\sum_jT_{lj}w_j+r_l(s,w),
\qquad \norm{r_l(s,w)}_2\le\eta_l .
\label{eq:affine}
\end{equation}
其中 $y_l^0$ 是参考潮流，$t_l$ 是新增负荷灵敏度，$T_{lj}$ 是光伏有功和无功灵敏度，$\eta_l$ 是二阶修正量。这样就可以同时处理有功和无功变化对线路电流的影响。

为计算最不利的复功率幅值，程序在 $0$ 至 $2\pi$ 范围内搜索角度 $\theta$。令 $e_\theta=(\cos\theta,\sin\theta)$，则仿射潮流的最大幅值可表示为
\begin{equation}
A_l=\max_{\theta\in[0,2\pi]}
\left[e_\theta^\top y_l^0+
h|e_\theta^\top t_l|+
\sum_j\supp_{\U_j}(T_{lj}^\top e_\theta)\right].
\label{eq:angular}
\end{equation}
该式只需要在一个角度变量上搜索；每个角度对应的光伏最不利影响由前述半圆支撑函数直接计算。设括号内的角度函数为 $f_l(\theta)$。实现中采用 $N_\theta$ 个等距角度，并按照
\begin{align}
K_l&=\norm{y_l^0}_2+h\norm{t_l}_2+
\sum_jC_j\norm{T_{lj}}_{\mathrm F},\label{eq:lip}\\
A_l&\le\max_{0\le k<N_\theta}
f_l(2\pi k/N_\theta)+\frac{\pi K_l}{N_\theta}
\label{eq:gridsafe}
\end{align}
补偿角度网格间的可能变化。结合线路接收端平方电压下界 $\underline\nu_l$ 和潮流二阶修正 $\eta_l$，得到该负荷单元内的电流平方计算界
\begin{equation}
\ell_l\le U_l^{\rm phasor}:=
\frac{\left[\max\left\{0,\max_k f_l(2\pi k/N_\theta)
+\frac{\pi K_l}{N_\theta}+\eta_l\right\}\right]^2}
{\underline\nu_l}.
\label{eq:phasor_bound}
\end{equation}
对每条线路比较 $U_l^{\rm phasor}$ 与 $(I_l^{\max})^2$，即可完成线路过载检查。正式容量算法只使用这一种线路电流计算界；分块 Hessian 主要用于节点电压和逆变器相关的非线性修正。

\subsection{逆变器视在功率检查}

光伏逆变器还需要满足自身的额定容量。对于第 $j$ 台逆变器，按实际有功出力与原始下垂无功命令形成
\begin{equation}
y_j^{PV}(s,w)=
\begin{pmatrix}
P_j^0+w_{p,j}\\
w_{q,j}-g_j(V_j(s,w)-V_j^{\rm ref})
\end{pmatrix}.
\end{equation}
与线路潮流类似，程序用一阶灵敏度计算其有功--无功向量的变化范围，再计入电压的二阶修正。对于电压幅值的修正，可由平方电压导数计算
\begin{equation}
\nabla^2 V_j=
\frac{\nabla^2\nu_j}{2\sqrt{\nu_j}}
-\frac{\nabla\nu_j\nabla\nu_j^\top}{4\nu_j^{3/2}}.
\end{equation}
记计算得到的仿射功率幅值上界为 $\overline A_j^{PV}$，非线性修正为 $\eta_j^{PV}$，对应的逆变器容量条件是
\begin{equation}
\overline A_j^{PV}+\eta_j^{PV}\le S_j,\qquad j=1,\ldots,m.
\label{eq:PVpass}
\end{equation}
完成节点电压、线路电流和逆变器视在功率三类检查后，当前负荷单元的安全结果传入下一节的容量搜索与协调增益更新。
